\chapter{Character counts from finite matrices}
\label{chap:library-computations}

Restriction of ordinary characters to regular elements gives a
matrix whose rank is the number of irreducible Brauer characters.
If an involution fixes the block, a second rank calculation gives
the number of Brauer characters it fixes. Both calculations follow
from linear algebra once the matrix entries and the action on
conjugacy classes are identified.

\section{The rank of ordinary restrictions}

\begin{lemma}
Let $M,N,V$ be vector spaces over a field $K$, with $N$ finite
dimensional. Let $D:M\to N$ and $E:N\to V$ be linear maps.
If $D$ is surjective and $E$ is injective, then
\[
 \dim_K\im(ED)=\dim_K N.
\]
\end{lemma}
\begin{proof}
Surjectivity gives $\im(ED)=\im E$. Injectivity identifies $\im E$
with $N$.
\end{proof}

We recall the character theoretic application over the complex
numbers, with a fixed choice of roots for Brauer characters.
For an $\ell$-block $b$, write $\Irr(b)$ and $\IBr(b)$ for its ordinary
irreducible characters and irreducible Brauer characters.

\begin{theorem}[Rank of ordinary restrictions]
Let $b$ be an $\ell$-block of a finite group $G$. If $\chi^0$ denotes
restriction to the $\ell$-regular elements, then
\begin{equation}
 \dim_{\mathbb C}\langle\chi^0:\chi\in\Irr(b)\rangle_{\mathbb C}
   =|\IBr(b)|.
 \label{lib:computations:restriction-rank}
\end{equation}
\end{theorem}
\begin{proof}
The irreducible Brauer characters are linearly independent.
The decomposition formula is
\[
 \chi^0=\sum_{\varphi\in\IBr(b)}d_{\chi\varphi}\varphi
 \qquad(\chi\in\Irr(b)).
\]
The decomposition matrix of $b$ has rank $|\IBr(b)|$.
Thus its rows span the coordinate space on $\IBr(b)$.
Apply the lemma with $D$ sending the ordinary character basis to
these rows and $E$ sending the coordinate basis to the Brauer
characters.
\end{proof}

For linear independence, the decomposition formula and full rank
of the decomposition matrix, see, for example,
\cite[Theorem~2.6 and Corollaries~2.9--2.11]{Navarro1998}.
The passage to one block follows from the block diagonal form of
the decomposition matrix
\cite[Theorem~3.3 and the following discussion]{Navarro1998}.
The related equality between $|\IBr(G)|$ and the number of regular
conjugacy classes is Corollary~2.10 in that account.

The same rank argument works over a field $K$ of characteristic
zero if the characters take values in $K$, the decomposition formula
holds there, the Brauer characters are linearly independent over
$K$, and the decomposition rows span the coordinate space over $K$.
Theorem~\ref{lib:brauer:independence} gives independence under its
stated field and root hypotheses. The decomposition and spanning
conditions are additional hypotheses.

\begin{implementation}
\module{Formalisation.RestrictionRank} proves
\lean{finrank_range_comp_eq_card} and
\lean{finrank_span_matrix_combinations_eq_card}.
The range calculation applies Mathlib's
\lean{LinearMap.range_comp_of_range_eq_top} and
\lean{LinearMap.finrank_range_of_inj}, followed by its dimension
formula for a finite coordinate space.
The second theorem assumes independent Brauer vectors and
decomposition rows spanning their coordinate space.
It proves the linear algebra implication over $K$.
The identification of these vectors and coefficients with characters
and decomposition numbers is separate.
\end{implementation}

Let $c_1,\ldots,c_t$ represent all $\ell$-regular conjugacy classes.
Evaluation at these representatives is injective on class functions
defined on the $\ell$-regular elements.
Consequently the matrix
\[
 M=(\chi(c_j))_{\chi\in\Irr(b),\,1\leq j\leq t}
\]
has row rank $|\IBr(b)|$. More generally, one may use any list
meeting every regular conjugacy class. To apply the formula to a
specified matrix, its rows must be the indicated ordinary characters
and its columns must evaluate those characters at such a list.

\section{Fixed points from symmetrised rows}

\begin{proposition}[An involution permuting a basis]
\label{lib:computations:fixed-rank}
Let $K$ have characteristic zero. Let $(v_i)_{i\in I}$ be a finite
basis of a $K$-vector space $V$, and let $T:V\to V$ be a linear map
such that $T(v_i)=v_{\sigma(i)}$ for an involution $\sigma$ of $I$.
Set
\[
 r=\dim_K V,\qquad s=\dim_K\im(1+T).
\]
Then $|I^\sigma|=2s-r$.
\end{proposition}
\begin{proof}
The permutation matrix of $T$ has diagonal entry one at a fixed
index and zero otherwise. Hence $\tr(T)=|I^\sigma|$ in $K$.
Since $T^2=1$, the map $P=(1+T)/2$ is a projection onto $\im(1+T)$.
The trace of a projection is its rank, so
\[
 s=\tr(P)=\tfrac12(r+\tr(T)).
\]
Characteristic zero gives the asserted equality of integers.
\end{proof}

The basis need not be known explicitly. Suppose $(g_a)_{a\in A}$
spans $V$ and $E:V\to K^C$ is an injective linear map.
Let $M$ be a matrix and let $\pi$ permute $C$, with
\[
 E(g_a)(c)=M_{a,c},\qquad
 E(Tg_a)(c)=M_{a,\pi(c)}.
\]
The vectors $(1+T)g_a$ span $\im(1+T)$. Their images under $E$ are
the rows of the \emph{symmetrised matrix}
\[
 M^+_{a,c}=M_{a,c}+M_{a,\pi(c)}.
\]
Thus $M$ has row rank $r$ and $M^+$ has row rank $s$.

For a block fixed by an involution in $\Aut(G)$, take the basis
to be its irreducible Brauer characters and the spanning vectors
to be its ordinary restrictions. Once $\pi$ describes the action
on the chosen regular class representatives, the two ranks give
\[
 \bigl(|\IBr(b)|,\ |\IBr(b)^\sigma|\bigr)=(r,2s-r).
\]
The action formula must use the same characters, roots and class
representatives as the evaluation formula.

\begin{implementation}
The trace and projection argument is
\lean{fixed_card_add_card_eq_two_mul_plus_rank} in
\module{ModularRep.PaperProofs.SporadicFi24P3Definition44NamedCarrierInvolutionBasisRank}.
It applies to any finite basis over a field of characteristic zero.
Its proof uses Mathlib's \lean{Matrix.trace_permutation} and
\lean{LinearMap.IsProj.trace}. The local argument identifies the
basis action with a permutation matrix and relates the projection
to the image of $1+T$.
The comparison with the symmetrised row rank is
\lean{finrank_range_id_add_eq_symmetrizedRows} in
\module{ModularRep.PaperProofs.SporadicFi24P3Definition44NamedCarrierSymmetrizedRowRankEquality}.
Its hypotheses are the spanning, injectivity and evaluation
identities above. The matrix identity in this latter theorem
allows any map on the column index set.
\end{implementation}

\section{Certifying a matrix rank}

\begin{proposition}[A minor and spanning rows]
Let $M$ be a finite matrix over a field $K$. Suppose a square minor
on $r$ selected rows and $r$ selected columns has nonzero determinant.
If every row of $M$ is a linear combination of the selected rows,
then $M$ has row rank $r$.
\end{proposition}
\begin{proof}
Restricting the selected rows to the selected columns gives the
rows of an invertible matrix, so they are linearly independent.
By hypothesis they span the row space of $M$. They therefore form
a basis of that space.
\end{proof}

\begin{implementation}
\module{ModularRep.PaperProofs.FiniteRowRankCertificate}
records the selected rows and columns, the determinant of the
minor, its nonvanishing and the coefficients expressing every row
as a combination of the selected rows.
\lean{FiniteRowRankCertificate.rows_finrank} proves the rank from
the determinant identity and the equations expressing every row
in terms of the selected rows.
It uses Mathlib's \lean{Matrix.linearIndependent_rows_of_det_ne_zero}
for the selected minor and \lean{finrank_span_eq_card} for the
dimension of the span. The local proof shows that the selected rows
span the full row space.
\end{implementation}

Together with the evaluation and action identities, these finite
equations give character counts and counts of fixed characters.
To use the rank calculation for character counts, one must identify the full
matrix with the restrictions of all the intended characters to all the relevant
regular classes. The selected minor proves that the selected rows are linearly
independent, and the given row expressions prove that they span the full row space.

\section{Finite scalar products and character selection}

Let $K$ be a field of characteristic zero, let $H$ be a finite group,
and let $(\mathcal C_j)_{j\in J}$ be distinct
conjugacy classes with representatives $u_j$, where $J$ is finite.
If a function $f:H\to K$ is constant on each $\mathcal C_j$ and vanishes
outside their union, partitioning the finite sum gives
\[
 \frac{1}{|H|}\sum_{g\in H}f(g)
   =\sum_{j\in J}\frac{|\mathcal C_j|}{|H|}f(u_j).
\]
Every coefficient $w_j=|\mathcal C_j|/|H|$ is nonzero. In a character calculation,
the required constancy and vanishing concern the product of the two
functions in the scalar product.

\begin{implementation}
\module{ManuscriptIBAW.Library.FiniteClassPairing} proves
\lean{FiniteClassPairing.sum_localize} for a finite set with an explicit
class map and distinct selected class labels. It uses constancy on
the selected fibres and vanishing elsewhere. The theorem
\lean{FiniteClassPairing.weight_ne_zero} proves nonvanishing of the
normalised fibre sizes in characteristic zero, using a representative
in each fibre. These results do not require a character theory.
\end{implementation}

Let $K$ be a field, let $J$ be finite, and let the vectors
$v_i\in K^J$, indexed by a set $I$, span $K^J$. Suppose that
$w_j\ne0$ for every $j$ and that $z_r\in K^J$ is nonzero for each
$r\in J$. Let $a_{i,r}\in K$ and $\varepsilon_i\in K$ satisfy
\[
 \sum_{j\in J}w_jv_i(j)z_r(j)=\varepsilon_i a_{i,r}
 \qquad(i\in I,\ r\in J).
\]
Then no column of $(a_{i,r})$ is zero. Indeed, if the column at $r$
were zero, the linear functional
$v\mapsto\sum_jw_jv(j)z_r(j)$ would vanish on a spanning family.
It would therefore vanish on every coordinate vector, giving
$w_jz_r(j)=0$ for all $j$. This contradicts $z_r\ne0$.

Suppose now that $K$ has characteristic zero, every $a_{i,r}$ is a
natural number viewed in $K$, and each row has sum one. Every row
then has exactly one entry equal to one and all others zero. Since
every column is nonzero, choose $s(r)\in I$ with $a_{s(r),r}=1$.
The map $s:J\to I$ is injective and satisfies
\[
 a_{s(r),j}=\begin{cases}1&j=r,\\0&j\ne r.\end{cases}
\]
Injectivity follows by comparing the unique nonzero entries in two
selected rows.

\begin{implementation}
In \module{ManuscriptIBAW.Library.FiniteCharacterSelection},
\lean{FiniteCharacterSelection.coverage_of_span} proves that every
column is nonzero. The theorem
\lean{FiniteCharacterSelection.exists_delta_injection_of_span}
constructs the injective selection assuming that the entries are
nonnegative integers and each row sums to one. These are finite linear algebra and arithmetic
results over the stated fields.
\lean{span_eq_top_of_diagonal_pairings} and
\lean{span_eq_top_of_trace_expansions} derive the spanning condition
from diagonal pairing identities with nonzero diagonal entries and from
explicit finite trace expansions, respectively.
\lean{linearIndependent_of_spanning_independent_family} proves that a finite family is linearly independent if its span contains
a linearly independent family of the same size.
In the character application, the identification of these families,
expansions and scalar products with the geometric and representation
theoretic results remains a separate set of assumptions.
\end{implementation}

The row sum condition can also be proved by counting. Let $I$ and $J$
be finite sets with $I\ne\varnothing$, and let $U$ be a set.
Let $f:J\to U$, $u\in U$, and $a:I\times J\to\mathbb N$. Suppose that the sum in column $j$ is
one when $f(j)=u$ and zero otherwise, and that $|I|=|f^{-1}(u)|$.
If every two rows are related by a permutation of $J$, then every row
has sum one. Indeed, all row sums have a common value $r$. Summing
the entries by rows and by columns gives
\[
 |I|r=|f^{-1}(u)|=|I|,
\]
and $I\ne\varnothing$ implies $r=1$.

\begin{implementation}
\lean{ManuscriptIBAW.TypeC.PrincipalGGGR.row_sum_one} proves this
finite arithmetic statement. It takes the column sums, the equality
of cardinalities and the row permutations as assumptions.
\end{implementation}
